% !TEX program = xelatex
% 2027 新领军学科能力二试 · 综合 · 社科（预测训练卷）
\documentclass[cjk=noto, seal=first]{mojing-exam}

\examtitle{2027 新领军学科能力测试 · 预测训练卷}
\secret{训练用卷 ★ 非官方试题}

\begin{document}

\begin{paper}[time={建议 1.5 小时（训练设置）}, full=100]{二试 · 综合 · 社科}

\begin{instructions}
  \item 官方只公布二试综合六选二的科目名称（社科与经济是两个独立科目），没有公布社科的考纲、题型与分值。本卷侧重社会科学中的证据与推理：
        集体决策、数据解读、抽样调查与因果判断，所需概念均在题中给出，全部安排为训练设置。
  \item 本卷共 4 题，每题 25 分，满分 100 分。
\end{instructions}

\begin{gradingnotes}
  \item 本评分标准为训练用建议标准。论述题按要点给分，言之成理、论证清楚者从高给分；分值为累计得分。
\end{gradingnotes}

\scoretable{1,2,3,4}

\begin{questions}

% =============================================================================
\question[25]
（集体如何做决定）三名委员甲、乙、丙对方案 A、B、C 的偏好排序如下（从最喜欢到最不喜欢）：甲：A、B、C；乙：B、C、A；丙：C、A、B。
\begin{enumerate}
  \item 用“两两多数表决”比较每一对方案，结果如何？“多数人的偏好”是否给出了一个合理的排序？
  \item 主席决定分两轮表决：先在两个方案中表决，胜者再与第三个方案表决。设各委员每轮都按真实偏好投票，说明主席只要安排表决顺序，就能使任何一个方案最终胜出；
  \item 若改用“波达计数”（每人给第一名 2 分、第二名 1 分、第三名 0 分，总分最高者胜，并列时抽签决定），结果如何？
  \item 某委员发现，如果他在投票时不按真实偏好排序，结果可能对他更有利。以 (2) 中的两轮表决为例说明这一点，并谈谈这类现象对设计集体决策规则的启示。
\end{enumerate}
\answerspace{2cm}
\begin{solution}
(1)~A 对 B：甲、丙支持 A，A 胜；B 对 C：甲、乙支持 B，B 胜；C 对 A：乙、丙支持 C，C 胜。出现循环 A$\succ$B$\succ$C$\succ$A，
多数偏好不具有传递性，得不到一致的排序（孔多塞悖论）。\score{8}

(2)~要让 A 胜：先 B 对 C（B 胜），再 B 对 A（A 胜）。同理，先表决另外两个方案，让想要的方案只参加最后一轮即可。\score{15}

(3)~三个方案各得 $2+1+0=3$ 分，三者并列，无法决出胜者。\score{19}

(4)~例：议程为“先 A 对 B，胜者再对 C”。都按真实偏好投票时，A 先胜 B，再败给 C，最终 C 胜出，这是甲最不喜欢的结果。
若甲在第一轮改投 B，则 B 胜 A，第二轮甲、乙都支持 B，B 胜 C，最终结果是甲的第二选择 B——甲不按真实偏好投票反而得益。\score{22}
不同的合理规则可以给出不同结果，结果可能取决于议程安排或策略性投票。设计规则时要权衡公平、稳定与抗操纵，并公开、预先确定程序；
可以提及：数学上已经证明，满足若干看似合理要求的“完美”排序规则不存在（阿罗不可能定理）。\score{25}
\end{solution}

% =============================================================================
\question[25]
（数据中的“悖论”）某校两个专业的录取数据如下（虚构）：
\begin{center}
\begin{tabular}{c|cc|cc}
  \toprule
  & \multicolumn{2}{c|}{男生} & \multicolumn{2}{c}{女生} \\
  & 申请 & 录取 & 申请 & 录取 \\
  \midrule
  专业 X & 100 & 80 & 20 & 18 \\
  专业 Y & 20 & 4 & 100 & 30 \\
  \bottomrule
\end{tabular}
\end{center}
\begin{enumerate}
  \item 分别计算两个专业中男、女生的录取率，以及全校男、女生的录取率；
  \item 有人根据全校数据说“学校歧视女生”，也有人根据分专业数据说“学校偏向女生”。解释为什么两组数字方向相反；
  \item 要比较录取是否对性别公平，应当怎样比较？还需要哪些信息才能下结论？
\end{enumerate}
\answerspace{2cm}
\begin{solution}
(1)~专业 X：男 80\%，女 90\%；专业 Y：男 20\%，女 30\%。全校：男 $\frac{84}{120}=70\%$，女 $\frac{48}{120}=40\%$。\score{8}

(2)~每个专业女生录取率都更高，但女生大多申请录取率很低的专业 Y，男生大多申请录取率高的专业 X。全校录取率是各专业录取率按申请人数加权的平均，
权重不同导致方向反转（辛普森悖论）。\score{17}

(3)~应在专业内部比较，或用相同的申请分布对两组加权（标准化）后比较。还需了解申请者的资质分布（同一专业中男女申请者条件是否可比）、
各专业的招生名额与评价标准等；仅凭录取率不能断定是否存在歧视。\score{25}
\end{solution}

% =============================================================================
\question[25]
（抽样调查）阅读材料，回答问题。

\hspace{2em}据史料记载，1936 年美国总统选举前，《文学文摘》杂志向约一千万人寄出模拟选票，名单主要来自电话簿、汽车登记和杂志订户，收回约二百多万份，
据此预测共和党候选人兰登将大胜。结果民主党候选人罗斯福以约六成的普选票获胜。同年，盖洛普用规模小得多（数万人）、按人口构成设计的样本，
正确预测了罗斯福获胜。（数字为约数。）
\begin{enumerate}
  \item 指出《文学文摘》调查中两种主要的偏差来源，并结合当时的社会背景（经济大萧条）说明它们为什么会使结果偏向兰登；
  \item 为什么“样本量巨大”没有能够弥补这些偏差？
  \item 如果由你设计一项调查，估计某城市高中生每天的平均睡眠时间，你会怎样抽样？说明如何减少类似的偏差。
\end{enumerate}
\answerspace{2cm}
\begin{solution}
(1)~选择偏差（抽样框偏差）：名单来自电话、汽车拥有者和杂志订户，在大萧条时期这些家庭中相对富裕、更可能支持共和党的人所占比例偏高。\score{6}
无回答偏差：只有约四分之一的人寄回选票，主动回复者的政治态度与不回复者可能系统地不同（例如对现任政府更不满的人更愿意表达）。
两种偏差都起作用；后来的研究对哪一种更主要有不同看法，答出两者并说明机制即可。\score{12}

(2)~样本量只能减小随机误差，不能消除系统偏差：从有偏的总体中抽得再多，估计的仍是那个有偏总体的特征，而且精度越高，错得越“自信”。\score{17}

(3)~例：取得全市高中学校与班级名单，按区域、学校类型分层，在各层内随机抽取班级，对抽中班级全体学生调查（整群抽样）；
匿名作答，提高回答率并追访未回答者；用睡眠日志而非单次回忆减少测量误差；按各层人数加权估计全市平均。
应采用概率抽样（每个学生被抽中的概率已知）：盖洛普当年用的配额抽样在 1948 年的选举预测中同样失败过。\score{25}
\end{solution}

% =============================================================================
\question[25]
（相关与因果）阅读下面两段论述，回答问题。

\hspace{2em}论述一：“统计显示，冰淇淋销量高的月份，溺水事故也多。因此，为减少溺水事故，应当限制夏季冰淇淋的销售。”

\hspace{2em}论述二：“某校参加课外数学小组的学生，数学成绩明显高于不参加的学生。因此，学校应要求所有学生参加数学小组，这样全校的数学成绩会显著提高。”
\begin{enumerate}
  \item 分别指出两段论述在推理上的问题；
  \item 设计一个能检验“参加数学小组能提高成绩”的研究方案，说明它为什么能排除 (1) 中指出的问题；
  \item 如果出于伦理或现实原因无法进行随机分组，你还能用什么办法使结论更可信？举一例说明。
\end{enumerate}
\answerspace{2cm}
\begin{solution}
(1)~论述一：两者都随气温升高而增加，存在共同原因（混杂因素），相关不等于因果。\score{5}
论述二：自愿参加者本来就可能更有兴趣、基础更好（自我选择），成绩差异未必由小组造成；即使有效，强制参加也可能因兴趣不同而效果不同。\score{11}

(2)~随机对照实验：在自愿报名但名额有限的学生中随机抽签决定谁参加，比较一段时间后两组成绩的变化（最好同时测量前测成绩）。
随机分组使两组在兴趣、基础等方面平均相同，差异只能归因于是否参加。\score{19}

(3)~例：比较同一批学生参加前后的成绩变化，并与同期未参加者的变化对比（双重差分）；或利用名额按报名顺序截止等“近似随机”的自然实验；
控制前测成绩、家庭背景等可观测变量，并说明仍可能存在未观测的混杂因素。\score{25}
\end{solution}

\end{questions}
\end{paper}
\end{document}
